\documentclass[11pt]{article}
\usepackage[margin=1in]{geometry}
\usepackage{enumitem}
\usepackage{amsmath,amssymb}
\usepackage{hyperref}

\setlist[enumerate]{itemsep=1.1em, topsep=0.6em}

\begin{document}

\begin{center}
{\LARGE \textbf{Probability 101 — Detailed \& Pedagogical Solution Sheet}}\\[0.2em]
{\small Solutions corresponding to \texttt{Probability\_101\_Brain\_Teasers.tex}}
\end{center}

\vspace{0.5em}
\noindent\textbf{Conventions.} Throughout, coins and dice are fair unless stated otherwise. For conditional probability, we use
\[
\mathbb{P}(A\mid B)=\frac{\mathbb{P}(A\cap B)}{\mathbb{P}(B)}.
\]

% =========================================================
\section*{Easy (7) — Solutions}

\begin{enumerate}[label=\textbf{\arabic*.}]

\item \textbf{Two uniforms: compute $\mathbb{P}(XY>1/2)$.}\\
Let $X,Y\sim \text{Unif}(0,1)$ independent. In the unit square $(x,y)\in[0,1]^2$ we want the area where $xy>1/2$.

For a fixed $x>0$, the condition is
\[
y>\frac{1}{2x}.
\]
If $x\le \tfrac12$, then $\tfrac{1}{2x}\ge 1$, so there are \emph{no} $y\in[0,1]$ satisfying it. Hence only $x\in(\tfrac12,1]$ contributes.

So
\[
\mathbb{P}(XY>1/2)=\int_{1/2}^1 \left(1-\frac{1}{2x}\right)\,dx
= \left[x-\frac12\ln x\right]_{1/2}^1.
\]
Compute:
\[
= \left(1-\frac12\ln 1\right) - \left(\frac12-\frac12\ln\frac12\right)
= 1-\frac12 + \frac12\ln\frac12
= \frac12-\frac12\ln 2.
\]
\[
\boxed{\mathbb{P}(XY>1/2)=\frac{1-\ln 2}{2}.}
\]
\hfill{\small\textit{Ref: Heard on the Street, Ch.\ 4, Q4.1}}

% ---------------------------------------------------------
\item \textbf{Coin flips: 4 coins vs 5 coins.}\\
Let $A\sim\text{Bin}(4,\tfrac12)$ be your number of heads, and $B\sim\text{Bin}(5,\tfrac12)$ your opponent’s heads, independent.
We want $\mathbb{P}(A>B)$.

A clean way is to sum over $b$:
\[
\mathbb{P}(A>B)=\sum_{b=0}^{5}\mathbb{P}(B=b)\,\mathbb{P}(A\ge b+1).
\]
Using binomial probabilities:
\[
\mathbb{P}(A=k)=\binom{4}{k}2^{-4},\qquad \mathbb{P}(B=b)=\binom{5}{b}2^{-5}.
\]
Compute $\mathbb{P}(A\ge b+1)$ for each $b$:
\[
\begin{array}{c|c}
b & \mathbb{P}(A\ge b+1)\\ \hline
0 & 1-\mathbb{P}(A=0)=1-\frac{1}{16}=\frac{15}{16}\\
1 & 1-\mathbb{P}(A\le 1)=1-\frac{1+4}{16}=\frac{11}{16}\\
2 & 1-\mathbb{P}(A\le 2)=1-\frac{1+4+6}{16}=\frac{5}{16}\\
3 & \mathbb{P}(A=4)=\frac{1}{16}\\
4 & 0\\
5 & 0
\end{array}
\]
Now sum:
\[
\mathbb{P}(A>B)=\frac{1}{32}\left[
\binom50\frac{15}{16}+\binom51\frac{11}{16}+\binom52\frac{5}{16}+\binom53\frac{1}{16}
\right].
\]
Factor $\frac{1}{32}\cdot\frac{1}{16}=\frac{1}{512}$:
\[
=\frac{1}{512}\left(1\cdot 15 + 5\cdot 11 + 10\cdot 5 + 10\cdot 1\right)
=\frac{1}{512}(15+55+50+10)=\frac{130}{512}=\frac{65}{256}.
\]
\[
\boxed{\mathbb{P}(\text{you win})=\frac{65}{256}.}
\]
\hfill{\small\textit{Ref: Heard on the Street, Ch.\ 4, Q4.3}}

% ---------------------------------------------------------
\item \textbf{Two kings in two dealt cards.}\\
There are $\binom{52}{2}$ equally likely 2-card hands. Favorable hands are choosing 2 kings from 4:
\[
\mathbb{P}(\text{both kings})=\frac{\binom{4}{2}}{\binom{52}{2}}
=\frac{6}{1326}=\frac{1}{221}.
\]
\[
\boxed{\frac{1}{221}.}
\]
\hfill{\small\textit{Ref: Heard on the Street, Ch.\ 4, Q4.12}}

% ---------------------------------------------------------
\item \textbf{Coin toss game: $A$ flips $n{+}1$, $B$ flips $n$. Find $\mathbb{P}(A>B)$.}\\
Write $A=A_n+Z$ where $A_n\sim\text{Bin}(n,\tfrac12)$ is $A$'s first $n$ coins and $Z\sim\text{Bern}(\tfrac12)$ is the extra coin (independent).
Let $B\sim\text{Bin}(n,\tfrac12)$ independent.

We condition on $Z$:
\[
\mathbb{P}(A>B)=\tfrac12\,\mathbb{P}(A_n>B)+\tfrac12\,\mathbb{P}(A_n+1>B).
\]
But $\mathbb{P}(A_n+1>B)=\mathbb{P}(A_n\ge B)$.

Now use symmetry: $A_n$ and $B$ are i.i.d., so
\[
\mathbb{P}(A_n>B)=\mathbb{P}(A_n<B).
\]
Let $p_0=\mathbb{P}(A_n=B)$ (tie probability). Then
\[
\mathbb{P}(A_n>B)=\frac{1-p_0}{2},
\qquad
\mathbb{P}(A_n\ge B)=1-\mathbb{P}(A_n<B)=1-\frac{1-p_0}{2}=\frac{1+p_0}{2}.
\]
Plug in:
\[
\mathbb{P}(A>B)=\tfrac12\cdot\frac{1-p_0}{2}+\tfrac12\cdot\frac{1+p_0}{2}
=\frac{1}{4}(1-p_0+1+p_0)=\frac12.
\]
\[
\boxed{\mathbb{P}(A>B)=\frac12 \text{ for all } n.}
\]
Pedagogical note: the ``extra coin'' helps exactly as much as the tie probability hurts, and the two effects cancel.

\hfill{\small\textit{Ref: A Practical Guide, Ch.\ 4, ``Coin toss game''}}

% ---------------------------------------------------------
\item \textbf{Card game (higher rank wins; tie loses).}\\
Draw your card, then dealer draws a card without replacement. By symmetry,
\[
\mathbb{P}(\text{your rank higher})=\mathbb{P}(\text{dealer rank higher}).
\]
So
\[
\mathbb{P}(\text{you win})=\frac{1-\mathbb{P}(\text{tie})}{2}.
\]
Compute $\mathbb{P}(\text{tie})$: tie means same rank. There are 13 ranks, each with 4 cards.
\[
\mathbb{P}(\text{tie})=\sum_{r=1}^{13}\frac{\binom{4}{2}}{\binom{52}{2}}
=13\cdot \frac{6}{1326}=\frac{78}{1326}=\frac{13}{221}.
\]
Therefore
\[
\mathbb{P}(\text{you win})=\frac{1-\frac{13}{221}}{2}
=\frac{\frac{208}{221}}{2}=\frac{104}{221}.
\]
\[
\boxed{\mathbb{P}(\text{win})=\frac{104}{221}.}
\]
\hfill{\small\textit{Ref: A Practical Guide, Ch.\ 4, ``Card game''}}

% ---------------------------------------------------------
\item \textbf{Three die rolls strictly increasing.}\\
Total outcomes: $6^3=216$.

Strictly increasing means $(a,b,c)$ with $1\le a<b<c\le 6$. Choose any 3 distinct faces from 6: $\binom{6}{3}=20$.
For each chosen set $\{a,b,c\}$ there is exactly \emph{one} increasing ordering.

So favorable outcomes $=20$, hence
\[
\mathbb{P}(\text{strictly increasing})=\frac{20}{216}=\frac{5}{54}.
\]
\[
\boxed{\frac{5}{54}.}
\]
\hfill{\small\textit{Ref: A Practical Guide, Ch.\ 4, ``Dice order''}}

% ---------------------------------------------------------
\item \textbf{Population standard deviation of $(1,2,3,4,5)$.}\\
Mean:
\[
\mu=\frac{1+2+3+4+5}{5}=3.
\]
Population variance:
\[
\sigma^2=\frac{1}{5}\sum_{i=1}^5 (x_i-\mu)^2
=\frac{1}{5}\big[(1-3)^2+(2-3)^2+(3-3)^2+(4-3)^2+(5-3)^2\big]
=\frac{1}{5}(4+1+0+1+4)=2.
\]
Standard deviation:
\[
\sigma=\sqrt{2}.
\]
\[
\boxed{\sigma=\sqrt{2}.}
\]
\hfill{\small\textit{Ref: Heard on the Street, Ch.\ 4, Q4.20}}

\end{enumerate}

% =========================================================
\section*{Medium (7) — Solutions}

\begin{enumerate}[label=\textbf{\arabic*.}, resume]

\item \textbf{Two jars (50 white, 50 black): maximize chance to draw a white.}\\
You place some whites/blacks in Jar 1 and the rest in Jar 2, then a jar is chosen uniformly and you draw uniformly from that jar.

Key intuition: you want \emph{one jar} to be as ``white-heavy'' as possible, because you get a 1/2 chance to sample from it.

A famous optimal strategy is:
\begin{itemize}
\item Jar 1: put \textbf{1 white} (and 0 black).
\item Jar 2: put the remaining 49 white and 50 black (99 total).
\end{itemize}
Then
\[
\mathbb{P}(\text{white})=\frac12\cdot 1 + \frac12\cdot \frac{49}{99}
=\frac12+\frac{49}{198}=\frac{74}{99}.
\]
\[
\boxed{\text{Max probability }=\frac{74}{99}\approx 0.7475.}
\]

Why is it optimal (sketch): if Jar 1 contains $w_1$ whites and $b_1$ blacks, your success probability is
\[
\frac12\frac{w_1}{w_1+b_1}+\frac12\frac{50-w_1}{100-(w_1+b_1)}.
\]
Adding a black to the ``special'' jar always hurts its purity; concentrating whites without blacks in one jar is best, and it is optimal to put exactly one white there (a standard calculus/convexity argument).

\hfill{\small\textit{Ref: Heard on the Street, Ch.\ 4, Q4.15}}

% ---------------------------------------------------------
\item \textbf{Two-headed coin posterior after 10 heads.}\\
Prior: choose the special two-headed coin with probability $1/1000$, otherwise a fair coin with probability $999/1000$.

Likelihood of 10 heads:
\[
\mathbb{P}(10H \mid \text{two-headed})=1,\qquad
\mathbb{P}(10H \mid \text{fair})=\left(\frac12\right)^{10}=\frac{1}{1024}.
\]
Bayes:
\[
\mathbb{P}(\text{two-headed}\mid 10H)
=\frac{\frac{1}{1000}\cdot 1}{\frac{1}{1000}\cdot 1 + \frac{999}{1000}\cdot \frac{1}{1024}}
=\frac{1}{1+\frac{999}{1024}}
=\frac{1024}{2023}.
\]
\[
\boxed{\mathbb{P}(\text{two-headed}\mid 10H)=\frac{1024}{2023}\approx 0.506.}
\]
Pedagogical note: despite 10 heads, the fair-coin explanation is still plausible because the fair-coin prior is huge (999x larger).

\hfill{\small\textit{Ref: Heard on the Street, Ch.\ 4, Q4.22}}

% ---------------------------------------------------------
\item \textbf{Medical test (base rates).}\\
Let $D$ be ``has disease''. Given:
\[
\mathbb{P}(D)=0.005,\quad
\mathbb{P}(+\mid D)=1,\quad
\mathbb{P}(+\mid D^c)=0.07.
\]
We want $\mathbb{P}(D\mid +)$:
\[
\mathbb{P}(D\mid +)=\frac{\mathbb{P}(+\mid D)\mathbb{P}(D)}{\mathbb{P}(+\mid D)\mathbb{P}(D)+\mathbb{P}(+\mid D^c)\mathbb{P}(D^c)}.
\]
Plug in:
\[
=\frac{1\cdot 0.005}{1\cdot 0.005 + 0.07\cdot 0.995}
=\frac{0.005}{0.005+0.06965}.
\]
As an exact fraction:
\[
0.005=\frac{5}{1000},\quad 0.07\cdot 0.995=\frac{7}{100}\cdot\frac{995}{1000}=\frac{6965}{100000},
\]
so
\[
\mathbb{P}(D\mid +)=\frac{100}{1493}\approx 0.06698.
\]
\[
\boxed{\mathbb{P}(D\mid +)=\frac{100}{1493}\approx 6.7\%.}
\]
Pedagogical note: false positives dominate because the disease is rare.

\hfill{\small\textit{Ref: Heard on the Street, Ch.\ 4, Q4.30}}

% ---------------------------------------------------------
\item \textbf{Drunk passenger (probability passenger 100 sits in seat 100).}\\
Classic invariant argument.

Passenger 1 picks a random seat. After that, every time someone finds their own seat taken, they choose randomly among the remaining seats.

Key observation: the \emph{only} seats that matter are Seat 1 and Seat 100.
Why? Because all seats 2--99, when freed, will be taken by their rightful owners as the process continues.

Eventually, some passenger will be forced to choose among the remaining seats, and the process ends when either:
\begin{itemize}
\item Seat 100 is taken by someone else (bad for passenger 100), or
\item Seat 1 is taken by someone else (then Seat 100 remains available for passenger 100).
\end{itemize}

At the moment the ``random choice'' happens with the two special seats still available, the chooser is equally likely to pick Seat 1 or Seat 100.

Therefore,
\[
\boxed{\mathbb{P}(\text{passenger 100 gets seat 100})=\frac12.}
\]
\hfill{\small\textit{Ref: A Practical Guide, Ch.\ 4, ``Drunk passenger''}}

% ---------------------------------------------------------
\item \textbf{Five letters in five envelopes: all wrong (derangement).}\\
We need the probability of a derangement of 5 items.

Total permutations: $5!=120$.

Let $A_i$ be event ``letter $i$ is in the correct envelope''.
We want $\mathbb{P}(\cap_{i=1}^5 A_i^c)$.

By inclusion--exclusion, number of derangements is:
\[
!5 = 5! - \binom51 4! + \binom52 3! - \binom53 2! + \binom54 1! - \binom55 0!.
\]
Compute:
\[
!5 = 120 - 5\cdot 24 + 10\cdot 6 - 10\cdot 2 + 5\cdot 1 - 1
=120-120+60-20+5-1=44.
\]
So
\[
\mathbb{P}(\text{all wrong})=\frac{44}{120}=\frac{11}{30}.
\]
\[
\boxed{\frac{11}{30}.}
\]
\hfill{\small\textit{Ref: A Practical Guide, Ch.\ 4, ``Application letters''}}

% ---------------------------------------------------------
\item \textbf{Two children: you meet a boy; probability both are boys.}\\
Interpretation consistent with the statement: pick a random two-child family (each child independently boy/girl with prob $1/2$), then you see \emph{one randomly selected child} from that family and it is a boy.

Possible family types (ordered): $BB, BG, GB, GG$, each with probability $1/4$.

Probability you meet a boy given each type:
\[
\mathbb{P}(\text{meet boy}\mid BB)=1,\quad
\mathbb{P}(\text{meet boy}\mid BG)=\frac12,\quad
\mathbb{P}(\text{meet boy}\mid GB)=\frac12,\quad
\mathbb{P}(\text{meet boy}\mid GG)=0.
\]
Bayes:
\[
\mathbb{P}(BB\mid \text{meet boy})
=\frac{\mathbb{P}(BB)\mathbb{P}(\text{meet boy}\mid BB)}{\sum \mathbb{P}(\text{type})\mathbb{P}(\text{meet boy}\mid \text{type})}
=\frac{\frac14\cdot 1}{\frac14\cdot 1+\frac14\cdot\frac12+\frac14\cdot\frac12}
=\frac{\frac14}{\frac12}=\frac12.
\]
\[
\boxed{\frac12.}
\]
Pedagogical note: this differs from the ``at least one boy'' version (which gives $1/3$) because the observation mechanism matters.

\hfill{\small\textit{Ref: A Practical Guide, Ch.\ 4, ``Boys and girls'' Part B}}

% ---------------------------------------------------------
\item \textbf{Monty Hall.}\\
Initial choice has probability $1/3$ of being the prize, $2/3$ of being a goat.

Host always opens a different door with a goat and never opens the prize door. Therefore:
\begin{itemize}
\item If your original choice was correct (prob $1/3$), switching loses.
\item If your original choice was wrong (prob $2/3$), switching wins (the only remaining unopened door must be the prize).
\end{itemize}

So switching wins with probability $2/3$:
\[
\boxed{\text{Switch; win probability }=\frac{2}{3}.}
\]
\hfill{\small\textit{Ref: A Practical Guide, Ch.\ 4, ``Monty Hall problem''}}

\end{enumerate}

% =========================================================
\section*{Hard (6) — Solutions}

\begin{enumerate}[label=\textbf{\arabic*.}, resume]

\item \textbf{Russian roulette with two adjacent bullets: spin or not?}\\
Model the cylinder as 6 chambers in a circle with pattern:
\[
B, B, E, E, E, E
\]
(two adjacent bullets, four empties). You spun once, pulled the trigger, heard a click, so your current chamber is empty.

\medskip
\noindent\textbf{Option (i): spin again.} Survival probability next shot is simply probability of landing on an empty chamber:
\[
\mathbb{P}(\text{survive}\mid \text{spin})=\frac{4}{6}=\frac{2}{3}.
\]

\medskip
\noindent\textbf{Option (ii): do not spin.} Given current chamber is empty, the next chamber is empty except in the single case where you are in the empty chamber immediately before the first bullet.

List the 4 empty chambers in the pattern above. Exactly 3 of them are followed by an empty chamber, and 1 is followed by a bullet. Hence
\[
\mathbb{P}(\text{survive}\mid \text{no spin, survived once})=\frac{3}{4}.
\]

\medskip
Compare:
\[
\frac{3}{4}=0.75 > \frac{2}{3}\approx 0.6667.
\]
\[
\boxed{\text{Do not spin; survival probability }=\frac{3}{4}.}
\]
\hfill{\small\textit{Ref: Heard on the Street, Ch.\ 4, Q4.21}}

% ---------------------------------------------------------
\item \textbf{Chocolate chips into 100 cookies: choose $N$ for $\mathbb{P}(\text{all nonempty})\ge 0.9$.}\\
Each chip independently lands in one of 100 cookies uniformly. Let $E_i$ be the event cookie $i$ is empty.

For a fixed cookie $i$:
\[
\mathbb{P}(E_i)=\left(1-\frac{1}{100}\right)^N=\left(\frac{99}{100}\right)^N.
\]
We want $\mathbb{P}(\cap_i E_i^c)\ge 0.9$, i.e.\ $\mathbb{P}(\cup_i E_i)\le 0.1$.

A simple \textbf{guarantee} uses the union bound:
\[
\mathbb{P}\Big(\bigcup_{i=1}^{100} E_i\Big)\le \sum_{i=1}^{100}\mathbb{P}(E_i)
=100\left(\frac{99}{100}\right)^N.
\]
So it suffices to choose $N$ such that
\[
100\left(\frac{99}{100}\right)^N \le 0.1
\quad\Longleftrightarrow\quad
\left(\frac{99}{100}\right)^N \le 0.001.
\]
Take logs:
\[
N\ln(0.99)\le \ln(0.001)=-6.907755\ldots
\]
Since $\ln(0.99)\approx -0.0100503$,
\[
N \ge \frac{6.907755}{0.0100503}\approx 687.2.
\]
Thus $N=688$ chips is sufficient:
\[
\boxed{N\approx 688 \text{ (guarantees } \ge 90\%\text{ by union bound).}}
\]
Pedagogical note: a sharper Poisson approximation gives about $686$, but $688$ is a safe clean answer.

\hfill{\small\textit{Ref: Heard on the Street, Ch.\ 4, Q4.32}}

% ---------------------------------------------------------
\item \textbf{Expected time to reach Box 4 (Markov chain).}\\
Let $E_i$ be the expected number of tosses to reach Box 4 starting from Box $i$.
Clearly $E_4=0$.

From Box 1: with probability $1/2$ go to 2, with probability $1/2$ go to 3, and you used one toss:
\[
E_1 = 1 + \frac12 E_2 + \frac12 E_3.
\]
From Box 2: with probability $1/2$ go to 1, with probability $1/2$ go to 4:
\[
E_2 = 1 + \frac12 E_1 + \frac12 E_4 = 1 + \frac12 E_1.
\]
From Box 3: Heads sends you to 2, Tails keeps you at 3:
\[
E_3 = 1 + \frac12 E_2 + \frac12 E_3.
\]
Solve the last one:
\[
E_3 - \frac12 E_3 = 1 + \frac12 E_2 \quad\Longrightarrow\quad
\frac12 E_3 = 1 + \frac12 E_2 \quad\Longrightarrow\quad
E_3 = 2 + E_2.
\]
Plug into $E_1$:
\[
E_1 = 1 + \frac12 E_2 + \frac12(2+E_2)=1+\frac12E_2+1+\frac12E_2=2+E_2.
\]
Now use $E_2=1+\tfrac12E_1$:
\[
E_2 = 1 + \frac12(2+E_2)=1+1+\frac12E_2
\quad\Longrightarrow\quad
\frac12E_2=2 \quad\Longrightarrow\quad E_2=4.
\]
Then
\[
E_1 = 2 + E_2 = 6.
\]
\[
\boxed{\mathbb{E}[\text{\# tosses from Box 1 to reach Box 4}]=6.}
\]
\hfill{\small\textit{Ref: Heard on the Street, Ch.\ 4, Q4.34}}

% ---------------------------------------------------------
\item \textbf{$N$ random points on a circle: probability all lie in some semicircle.}\\
This is a classic result:
\[
\boxed{\mathbb{P}(\text{all $N$ points lie in some semicircle})=\frac{N}{2^{N-1}}.}
\]

\medskip
\noindent\textbf{Proof idea (rotation + ``leftmost point'' argument).}
Represent points by angles on $[0,2\pi)$.
With probability 1, there is a unique point with the smallest angle; call it point $i$.

If all points lie in \emph{some} semicircle, then in particular they lie in the semicircle of length $\pi$ that starts at the smallest-angle point and goes counterclockwise. (If a semicircle covers all points, slide its left endpoint until it hits the smallest point.)

So, for a \emph{fixed} point $i$ to be the left endpoint witness, the other $N-1$ points must all fall in the next half of the circle, which has probability $(1/2)^{N-1}$.

These events for different $i$ are disjoint (only one point can be the unique smallest-angle point). Therefore,
\[
\mathbb{P}(\text{all in some semicircle})
= N\cdot \left(\frac12\right)^{N-1}
=\frac{N}{2^{N-1}}.
\]
\hfill{\small\textit{Ref: A Practical Guide, Ch.\ 4, ``$N$ points on a circle''}}

% ---------------------------------------------------------
\item \textbf{Single-elimination tournament: probability players 1 and 2 meet in the final.}\\
There are $2^n$ players. Pairings each round are random among remaining players. Better player always wins.

Players 1 and 2 meet in the final \emph{iff} they start in opposite halves of the bracket (so they cannot meet before the final).

Fix player 1's initial position. Among the remaining $2^n-1$ positions, exactly $2^{n-1}$ positions are in the opposite half of the bracket.

Thus
\[
\mathbb{P}(\text{1 and 2 meet in final})
=\frac{2^{n-1}}{2^n-1}.
\]
\[
\boxed{\frac{2^{n-1}}{2^n-1}.}
\]
Sanity checks: for $n=1$ (2 players) probability $=1$; for $n=2$ (4 players) probability $=2/3$.

\hfill{\small\textit{Ref: A Practical Guide, Ch.\ 4, ``Chess tournament''}}

% ---------------------------------------------------------
\item \textbf{Pirates, locks, keys (11 pirates; any 6 can open, no 5 can).}\\
We want the \emph{minimum} number of locks and keys per pirate.

\medskip
\noindent\textbf{Key observation.}
To ensure \emph{no group of 5} can open the safe, for each 5-person group $S$ there must exist \emph{at least one lock} for which \emph{none} of the pirates in $S$ has the key.

That means: for that lock, the set of pirates \emph{without} keys must include $S$.

But we also require \emph{any 6} pirates can open the safe. If a lock had 6 (or more) pirates without keys, that particular group of 6 would fail to open it. Therefore, for every lock:
\[
\text{number of pirates without that key} \le 5.
\]
Combining both constraints, any ``blocking'' set must have size exactly 5. Hence:
\begin{itemize}
\item Every lock corresponds to \emph{exactly one} 5-person group who does \emph{not} get the key.
\item To block \emph{every} 5-person group, we need at least one lock per 5-person group.
\end{itemize}

Number of 5-person groups among 11 pirates:
\[
\binom{11}{5}=462.
\]
So at least 462 locks are necessary.

\medskip
\noindent\textbf{Construction achieving 462 locks (so it is optimal).}
For each 5-person group $S\subset\{1,\dots,11\}$ with $|S|=5$, create one lock $L_S$ such that:
\begin{itemize}
\item Pirates in $S$ \emph{do not} receive a key to $L_S$.
\item The other 6 pirates (the complement of $S$) \emph{do} receive keys to $L_S$.
\end{itemize}
Now:
\begin{itemize}
\item Any group of 5 pirates equals some $S$, and they cannot open lock $L_S$, so they cannot open the safe.
\item Any group of 6 pirates is never contained in any 5-set $S$, so for every lock $L_S$ at least one of the 6 pirates has the key. Hence they can open \emph{all} locks, thus the safe.
\end{itemize}
Therefore, 462 locks suffice and are minimal:
\[
\boxed{\text{Minimum locks}=\binom{11}{5}=462.}
\]

\medskip
\noindent\textbf{How many keys does each pirate carry?}
A pirate gets a key for lock $L_S$ exactly when they are \emph{not} in the excluded set $S$.
Fix a pirate $i$. The number of 5-sets $S$ that \emph{do not contain} $i$ is:
\[
\binom{10}{5}=252
\]
(choose 5 excluded pirates from the other 10).
So each pirate carries:
\[
\boxed{\binom{10}{5}=252 \text{ keys.}}
\]

\hfill{\small\textit{Ref: A Practical Guide, Ch.\ 4, ``Screwy pirates''}}

\end{enumerate}

\end{document}